\scanpage{83}
Soit $\Pi_0$ le groupe engendré par les générateurs $\underline{\unsure{h|\ell}}_0,\underline h_1,\underline h_2$, soumis à

\[
 \underline h_0\,\underline h_1\,\underline h_2=1
\]

et soit $\deleted{\unsure{\widehat{\vphantom\Pi}}}\widehat\Pi$ son complété profini. Un hom de $\Pi_0$ dans un groupe $H$ (ou de $\widehat\Pi$ ds un gpe $H$ fini) équivaut à la donnée de

\[
 h_0,h_1,h_2\in H\qquad\text{satisfaisant}\quad h_0h_1h_2=1
\]

\margin{Soit $\mathcal T(H)\inserted{\subset H^3}$ l'ens des ces triples}

L'hom $\Pi_0\text{---}H$ ($\widehat\Pi\unsure{\text{---}|\longrightarrow}H$) est surjectif ssi $(h_0,h_1,h_2)$ est un syst. de \unsure{Schwartz|Schwarz}

\margin{$\operatorname{Sch}(H)\subset\mathcal T(H)$}

Soit \deleted{$\Pi$ le groupe des}

\diagram{\begin{tikzcd}[ampersand replacement=\&,column sep=tiny,row sep=small]
 \Omega_0 \arrow[d] \& {}=\operatorname{Aut}(\Pi_0)\\
 \Omega \& {}=\operatorname{Aut}(\widehat\Pi)
\end{tikzcd}}

Alors $\Omega_0$ (resp. $\Omega$, si $H$ fini) opère \inserted{à droite} sur \inserted{$\mathcal T(H)\simeq$} $\underline{\operatorname{Hom}}(\Omega,H)$ (\deleted{\unsure{sur}} $\simeq\underline{\operatorname{Hom}}(\widehat\Omega,H)$ \inserted{si $H$ fini}) et sur le sous-ensemble $\operatorname{Schw}(H)$. On a

\diagram{\begin{tikzcd}[ampersand replacement=\&,row sep=small,column sep=small]
 \& \& \& \operatorname{Autext}(\Pi_0)\arrow[d,equal] \&\\
 1\arrow[r,dash] \& \deleted{\unsure{\Pi}}\unsure{\Omega|\Pi}_0\arrow[r,dash] \& \Omega_0\arrow[r] \& \widetilde\Omega_0\arrow[r] \& 1\\
 1\arrow[r,"?",dash] \& \widehat\Pi\arrow[r,dash] \& \Omega\arrow[r] \& \widetilde\Omega\arrow[r,dash]\arrow[d,equal] \& 1\\
 \& \& \& \operatorname{Autext}(\widehat\Omega) \&
\end{tikzcd}}

et deux éléments de \deleted{$\underline{\operatorname{Hom}}(\Omega,H)$}\inserted{$\mathcal T(H)$} conjugués sous $\Pi$ ($\widehat\Pi$) sont conjugués dans $H$ -- l'inverse étant vrai pour les éléments de $\operatorname{Schw}(H)$. Ainsi $\widetilde\Omega_0$ ($\widetilde\Omega$) opère sur $\widetilde{\mathcal T}(H)=\mathcal T(H)/H$ et sur $\widetilde{\operatorname{Sch}}(H)=\operatorname{Sch}(H)/H=\operatorname{Sch}(H)/\underset{\text{ou }\widehat\Pi}{\Pi_0}$.

\deleted{\unsure{Dans le cas}} Dans une énumération des syst. de Schwartz relatifs à $H$, il convient de tenir compte des opérations (\unsure{ou} au \uline{moins} \unsure{quelques unes|quelques-unes}) de $\Omega_0$, ou $\Omega$ -- et de de $=$ pour $\widetilde\Omega_0,\widetilde\Omega$ quand on envisage $\widetilde{\operatorname{Sch}}(H)$.

\deleted{\unsure{$0\longrightarrow S_3\longrightarrow{}$} \unsure{$1\to\Omega\to S_3$}}
% archive sheet number: 82
